\scanpage{133}
\deleted{\unsure{(i)}}

$S$ \unsure{préschém.|préschéma} loc \unsure{noeth.}

$F$ foncteur \quad $(\mathrm{Sch})_{/S}\longrightarrow(\mathrm{Ens})$

(i) $F$ compatible avec limites inductives d'anneaux

(ii) $F$ compatible avec descente fidèlement plate \unsure{quasi cpcte.|quasi-compactes} \inserted{et localisation}

$\bigl[$(iii) Si $A$ anneau local noeth. \uline{complet} sur $S$,
corps résiduel $k$, alors $F(A)\overset{\sim}{\longrightarrow}F(k)$,

(iv) $F$ séparé.

(v) Pour \unsure{tt} corps $k$ sur $S$, $F(k)$ est \uline{fini},
et si $\unsure{K'|k'}$ est une ext. alg. close de $k$, \deleted{\unsure{$F(k$}} alg. clos,
$F(k)\unsure{-|\to}F(k')$ est bijectif.

Alors $F$ représentable par un \deleted{foncteur} objet
étale \deleted{\unsure{et qu.|et gpe}} séparé \unsure{qu.|quasi} fini sur $S$.

\deleted{\unsure{Réduction au cas|Réduisons-nous|Démonstration.}}

Supposons $S$ le spectre d'un anneau local complet
$A$, \deleted{et que} \inserted{à corps résiduel $k$ alg. clos.} de dim \unsure{$n$}, et qu'on sache
que sur $S_1=S-y$, le \deleted{foncteur} $F$ soit représentable
par $M_1$. \deleted{\unsure{Considérons les \illegible{1}}} Soient
$\unsure{\xi|\zeta}_i^0$ \deleted{\unsure{$\in F$|$\in k$}} \inserted{$(i\in I)$} les éléments de $\deleted{\unsure{R}}F(k)$, ils proviennent d'éléments
$\xi_i\in F(A)$ bien déterminés, considérons leurs
restrictions $\xi'_i=\xi_i\mid S_1\in F(S_1)$, ils correspondent à des
% archive sheet number: 132
